\chapter{Dualidad de \texorpdfstring{$L^p$ para $1<p<\infty$}{Lp para 1<p<infinito}}\label{v2:cap10}\label{cap:v2-10}
\section{De Hölder a la representación de funcionales}\label{lectura:v2-20}
Hölder construye un funcional continuo a partir de una función integrable. La pregunta inversa es más profunda: ¿todo funcional tiene ese origen? Las lecturas 20 y 21 de \cite{valencia2} responden afirmativamente mediante la geometría de $L^p$, sin imponer finitud ni $\sigma$-finitud a la medida. Mantendremos el emparejamiento bilineal $\int fg$; por eso las conjugaciones aparecerán en las funciones que realizan la norma y no en la definición del funcional. Esta convención hace lineal la identificación del dual.

\begin{proposicion}\label{v2:dual-isometria}
Sean $1<p<\infty$ y $q=p/(p-1)$. Si $h\in L^q$, entonces $\Phi(h)(f)=\int_Xhf\,d\mu$ define un funcional lineal continuo sobre $L^p$ y $\norm{\Phi(h)}=\norm h_q$.
\end{proposicion}
\begin{proof}
Hölder prueba $|\Phi(h)(f)|\le\norm h_q\norm f_p$. Para la desigualdad contraria, si $h\ne0$ en $L^q$, ponemos $f=\overline h|h|^{q-2}$ donde $h\ne0$ y $f=0$ donde $h=0$. Como $(q-1)p=q$, tenemos $|f|^p=|h|^q$ y
\[
 \norm f_p=\norm h_q^{q/p}=\norm h_q^{q-1},\qquad
 \Phi(h)(f)=\int|h|^q=\norm h_q^q.
\]
El cociente entre el valor del funcional y la norma de $f$ es $\norm h_q$. Si $h=0$, ambas normas son cero. La definición por separado en los ceros evita escribir una potencia negativa de cero.
\end{proof}

\section{Las desigualdades escalares de Clarkson}
\begin{lema}\label{v2:clarkson-escalar-mayor}
Si $p\ge2$, entonces
\[
 \left(\frac{1+x}{2}\right)^p+\left(\frac{1-x}{2}\right)^p\le\frac{1+x^p}{2}\quad(0\le x\le1),
\]
y, para $z,w\in\C$,
\[
 \left|\frac{z+w}{2}\right|^p+\left|\frac{z-w}{2}\right|^p
 \le\frac{|z|^p+|w|^p}{2}.
\]
\end{lema}
\begin{proof}
Para $0<x\le1$, sea
\[
 F(x)=x^{-p}\bigl((1+x)^p+(1-x)^p-2^{p-1}(1+x^p)\bigr).
\]
La regla del producto da
\[
 F'(x)=\frac p{x^{p+1}}\bigl(2^{p-1}-(1+x)^{p-1}-(1-x)^{p-1}\bigr).
\]
La suma $(1+x)^{p-1}+(1-x)^{p-1}$ es creciente, pues su derivada es $(p-1)((1+x)^{p-2}-(1-x)^{p-2})\ge0$, y en $x=1$ vale $2^{p-1}$. Por tanto $F'\ge0$, y $F(x)\le F(1)=0$. El extremo $x=0$ se verifica directamente. Dividir entre $2^p$ da la primera desigualdad.

Si $z=0$ o $w=0$, la segunda se reduce a $2^{1-p}\le1/2$. En otro caso, por simetría suponemos $|w|\le|z|$ y escribimos $w/z=re^{i\theta}$, $0<r\le1$. La expresión que hay que maximizar, salvo el factor $|z|^p/2^p$, es
\[
 H(c)=(1+r^2+2rc)^{p/2}+(1+r^2-2rc)^{p/2},\qquad c=\cos\theta.
\]
Es par en $c$ y creciente para $0\le c\le1$, porque
$H'(c)=pr((1+r^2+2rc)^{p/2-1}-(1+r^2-2rc)^{p/2-1})\ge0$.
Su máximo está en $|c|=1$, donde se aplica la primera desigualdad con $x=r$. Los exponentes $p$ en los dos sumandos y el signo menos entre las potencias de la derivada son necesarios.
\end{proof}
\begin{lema}\label{v2:clarkson-escalar-menor}
Si $1<p\le2$ y $q=p/(p-1)$, entonces
\[
 (1+x)^q+(1-x)^q\le2(1+x^p)^{q/p}\quad(0\le x\le1),
\]
y $|z+w|^q+|z-w|^q\le2(|z|^p+|w|^p)^{q/p}$ para $z,w\in\C$.
\end{lema}
\begin{proof}
Para $p=2$ hay igualdad. Supongamos $1<p<2$. Los extremos $x=0,1$ se verifican directamente. Sustituyendo $x=(1-u)/(1+u)$, $0<u<1$, y elevando a $p-1$, la primera desigualdad equivale a
\[
 (1+u^q)^{p-1}\le\tfrac12((1+u)^p+(1-u)^p).
\]
Las series binomiales convergen absolutamente para $0<u<1$. Al restarlas y agrupar el término de grado par de la primera serie con dos términos consecutivos de la segunda, la diferencia es $\sum_{k\ge1}T_k$, donde
\[
 T_k=\binom p{2k}u^{2k}-\binom{p-1}{2k-1}u^{q(2k-1)}-\binom{p-1}{2k}u^{2kq}.
\]
La fórmula $\binom ak=a(a-1)\cdots(a-k+1)/k!$ permite factorizar
\[
 T_k=\frac{u^{2k}}{(2k-1)!}\prod_{j=2}^{2k}(j-p)
 \left(\frac{1-u^{a_k}}{a_k}-\frac{1-u^{b_k}}{b_k}\right),
 \quad a_k=\frac{2k-p}{p-1},\quad b_k=\frac{2k}{p-1}.
\]
Para comprobar la factorización se expanden los dos cocientes: el término constante es $p(p-1)/(2k(2k-p))$, y los otros coeficientes son $-(p-1)/(2k-p)$ y $(p-1)/(2k)$. Son exactamente los tres coeficientes binomiales después de extraer el producto escrito. Los factores exteriores son positivos. Además $0<a_k<b_k$ y
\[
 \frac{1-u^a}{a}=\int_u^1t^{a-1}\,dt
\]
es decreciente en $a>0$, porque $0<t<1$. Así $T_k\ge0$, y la suma es no negativa. No se requiere afirmar convergencia uniforme de estas series en los extremos.

Para números complejos se reduce por homogeneidad a $z=1$, $w=re^{i\theta}$, $0\le r\le1$. Como $q\ge2$, el argumento de $H(c)$ en el lema anterior, ahora con exponente $q$, prueba que el máximo se alcanza cuando los dos números están alineados. La primera desigualdad del presente lema concluye.
\end{proof}
\begin{lema}[Desigualdad inversa para funciones no negativas]\label{v2:minkowski-inversa}
Si $0<r\le1$ y $a,b\ge0$ son medibles con integrales de sus potencias finitas, entonces
\[
 \left(\int(a+b)^r\right)^{1/r}\ge
 \left(\int a^r\right)^{1/r}+\left(\int b^r\right)^{1/r}.
\]
\end{lema}
\begin{proof}
Llamemos $A,B$ a los dos términos de la derecha. Si alguno es cero, su función es cero casi en todas partes. En otro caso, por concavidad de $t^r$,
\[
 \left(\frac{a+b}{A+B}\right)^r\ge
 \frac A{A+B}\left(\frac aA\right)^r+
 \frac B{A+B}\left(\frac bB\right)^r.
\]
Al integrar, el lado derecho vale uno. Tomar raíces prueba la afirmación. Esta es una desigualdad inversa para funciones no negativas; no es la desigualdad triangular de una norma cuando $r<1$.
\end{proof}
\begin{teorema}[Clarkson]\label{v2:clarkson}\index{Clarkson, desigualdades de}
Para $f,g\in L^p$, se cumplen
\begin{align*}
 p\ge2:&\quad\left\|\frac{f+g}{2}\right\|_p^p+\left\|\frac{f-g}{2}\right\|_p^p
 \le\frac{\norm f_p^p+\norm g_p^p}{2},\\
 1<p\le2:&\quad\left\|\frac{f+g}{2}\right\|_p^q+\left\|\frac{f-g}{2}\right\|_p^q
 \le\left(\frac{\norm f_p^p+\norm g_p^p}{2}\right)^{q/p}.
\end{align*}
\end{teorema}
\begin{proof}
La primera desigualdad se obtiene integrando el lema \ref{v2:clarkson-escalar-mayor}. Para la segunda, ponemos $r=p-1$, $u=|(f+g)/2|^q$ y $v=|(f-g)/2|^q$. Como $qr=p$, el lema \ref{v2:minkowski-inversa} da
\[
 \left\|\frac{f+g}{2}\right\|_p^q+\left\|\frac{f-g}{2}\right\|_p^q
 \le\left(\int(u+v)^{p-1}\,d\mu\right)^{1/(p-1)}.
\]
El lema escalar menor implica $u+v\le2^{1-q}(|f|^p+|g|^p)^{1/(p-1)}$. Elevando a $p-1$, el factor es $2^{(1-q)(p-1)}=1/2$. Integrar produce la cota anunciada.
\end{proof}

\section{Un funcional alcanza su norma}\label{lectura:v2-21}
\begin{teorema}\label{v2:lp-alcanza}
Todo funcional continuo no nulo $\varphi\in(L^p)^*$, $1<p<\infty$, alcanza su norma: existe $v$ con $\norm v_p=1$ y $\varphi(v)=\norm\varphi$.
\end{teorema}
\begin{proof}
Elegimos $u_n$ de norma uno con $|\varphi(u_n)|\to\norm\varphi$ y valores no nulos. Multiplicando $u_n$ por $\overline{\varphi(u_n)}/|\varphi(u_n)|$ obtenemos $v_n$ de norma uno con $\varphi(v_n)=|\varphi(u_n)|>0$. Es esta fase conjugada la que elimina el argumento complejo.

Si $(v_n)$ no fuera de Cauchy, existirían $\eps>0$ y pares de índices $a_j,b_j\to\infty$ con $\norm{v_{a_j}-v_{b_j}}_p\ge\eps$. Clarkson, con $s=p$ si $p\ge2$ y $s=q$ si $p\le2$, implica
\[
 \left\|\frac{v_{a_j}+v_{b_j}}2\right\|_p\le
 \left(1-(\eps/2)^s\right)^{1/s}=c<1.
\]
Pero $\varphi((v_{a_j}+v_{b_j})/2)\to\norm\varphi$, mientras que su valor absoluto es a lo sumo $c\norm\varphi$. Es una contradicción. Por completitud, $v_n\to v$; la continuidad de la norma y de $\varphi$ da las dos igualdades requeridas.
\end{proof}
\begin{lema}[McShane]\label{v2:mcshane}
Sean $V$ normado complejo y $\varphi\ne0$ continuo, con $\norm v=1$ y $\varphi(v)=\norm\varphi$. Si $\psi_w(t)=\norm{v+tw}$ y $\psi_{-iw}$ son diferenciables en $0$, entonces
\[
 \frac{\varphi(w)}{\norm\varphi}=\psi_w'(0)+i\psi_{-iw}'(0).
\]
En el caso real basta la primera derivada.
\end{lema}
\begin{proof}
Normalizamos $\norm\varphi=1$. Para todo $t\in\R$, $\norm{v+tw}\ge\operatorname{Re}\varphi(v+tw)=1+t\operatorname{Re}\varphi(w)$. Restar uno y dividir por $t$ da una cota inferior si $t>0$ y superior si $t<0$. La existencia de la derivada obliga a $\psi_w'(0)=\operatorname{Re}\varphi(w)$. Sustituir $w$ por $-iw$ da $\psi_{-iw}'(0)=\operatorname{Im}\varphi(w)$. Esta prueba evita una factorización con denominadores complejos innecesaria para el argumento.
\end{proof}
\begin{teorema}[Representación y reflexividad de $L^p$]\label{v2:dual-lp}\index{reflexividad!de Lp@de $L^p$}
Para cualquier espacio de medida y $1<p<\infty$, la aplicación $\Phi:L^q\to(L^p)^*$, $\Phi(h)(f)=\int hf\,d\mu$, es una isometría lineal sobreyectiva. En consecuencia, $L^p$ es reflexivo.
\end{teorema}
\begin{proof}
Ya se probó que es una isometría. Sea $\varphi\ne0$ y elijamos $v$ como en el teorema \ref{v2:lp-alcanza}. Para $w\in L^p$, estudiamos $I(t)=\int|v+tw|^p$. La función escalar $t\mapsto|a+tb|^p$ tiene derivada
\[
 p|a+tb|^{p-2}\operatorname{Re}(\overline{a+tb}\,b),
\]
interpretada como cero cuando $a+tb=0$. En efecto, allí el cociente incremental es $|b|^p|h|^p/h\to0$ porque $p>1$. Para $0<|t|\le1$, el teorema del valor medio da la cota
\[
 \left|\frac{|v+tw|^p-|v|^p}{t}\right|
 \le p(|v|+|w|)^{p-1}|w|.
\]
El lado derecho es integrable por Hölder, pues $(|v|+|w|)^{p-1}\in L^q$. Aplicamos convergencia dominada a los cocientes incrementales mismos, que son medibles; no hace falta elegir mediblemente un punto intermedio del teorema del valor medio. Así
\[
 I'(0)=p\int|v|^{p-2}\operatorname{Re}(\overline v w)\,d\mu.
\]
Como $I(0)=1$, la regla de la cadena da
$\psi_w'(0)=\int|v|^{p-2}\operatorname{Re}(\overline v w)\,d\mu$.
McShane recupera las partes real e imaginaria y obtiene
\[
 \varphi(w)=\norm\varphi\int |v|^{p-2}\overline v w\,d\mu.
\]
La función $h=\norm\varphi\,|v|^{p-2}\overline v$, definida como cero en los ceros de $v$, pertenece a $L^q$ porque $(p-1)q=p$. Representa a $\varphi$; el funcional cero se representa por $h=0$.

Finalmente, si $\Lambda\in(L^p)^{**}$, el funcional $h\mapsto\Lambda(\Phi(h))$ sobre $L^q$ está representado por un $f\in L^p$, aplicando el resultado con exponente $q$. Por tanto $\Lambda(\Phi(h))=\int fh=\Phi(h)(f)$ para todo $h$. Como todos los funcionales de $(L^p)^*$ son de la forma $\Phi(h)$, se sigue $\Lambda=J(f)$, donde $J$ es el encaje canónico en el bidual. Esta es la reflexividad.
\end{proof}

\section{Ejercicios resueltos: geometría, multiplicadores y compacidad}
\begin{ejercicio}[Lectura 21, ejercicio 1: precisión de la igualdad]\label{ej:v2-l21-1}
Para $1<p<\infty$ y $f,g\in L^p$ de normas no nulas, demostrar
\[
 2^{-|2-p|/p}\le
 \frac{\norm{f+g}_p^2+\norm{f-g}_p^2}{2(\norm f_p^2+\norm g_p^2)}
 \le2^{|2-p|/p}.
\]
Encontrar pares que alcancen ambos extremos cuando existen dos conjuntos medibles \emph{disjuntos} de medida positiva finita. La disjunción falta en la condición original para esta última petición.
\end{ejercicio}
\begin{solucion}
Escribamos $r=p/2$ y, para $h\ge0$, $Q_r(h)=(\int h^r)^{1/r}$. Si $r\ge1$, Minkowski y $a^r+b^r\le(a+b)^r$ dan
\[
 Q_r(a+b)\le Q_r(a)+Q_r(b)\le2^{1-1/r}Q_r(a+b).
\]
La segunda cota usa primero la desigualdad entre normas de un par de números: $A+B\le2^{1-1/r}(A^r+B^r)^{1/r}$. Si $0<r\le1$, la desigualdad inversa y $a^r+b^r\le2^{1-r}(a+b)^r$ dan
\[
 Q_r(a)+Q_r(b)\le Q_r(a+b)\le2^{1/r-1}(Q_r(a)+Q_r(b)).
\]
Para la última desigualdad se usa $(a+b)^r\le a^r+b^r$ y la convexidad de $t^{1/r}$ sobre los dos números $Q_r(a)^r,Q_r(b)^r$.

Aplicamos estas comparaciones a $a=|f+g|^2$, $b=|f-g|^2$ y usamos la identidad escalar $a+b=2(|f|^2+|g|^2)$. Resulta
\[
 \norm{f+g}_p^2+\norm{f-g}_p^2
 \le2C(\norm f_p^2+\norm g_p^2),\qquad C=2^{|1-2/p|}.
\]
Sustituyendo ahora $(f,g)$ por $(f+g,f-g)$ en esta cota se obtiene la desigualdad inferior, pues las nuevas suma y diferencia son $2f,2g$.

Sean $u,v$ indicadores normalizados de dos conjuntos disjuntos de medida positiva finita. Para $(f,g)=(u,v)$ el cociente es $2^{2/p-1}$. Para $(f,g)=(u+v,u-v)$ es $2^{1-2/p}$. Son los dos extremos, cuyo orden depende de si $p<2$ o $p>2$. Si $p=2$, ambos valen uno.
\end{solucion}
\begin{ejercicio}[Lectura 21, ejercicio 2]\label{ej:v2-l21-2}
Sea $f:X\to\F$ medible. Supóngase que $\{f\ne0\}$ es de medida $\sigma$-finita y que $fg\in L^1$ para todo $g\in L^p$, $1<p<\infty$. Demostrar $f\in L^q$, $q=p/(p-1)$.
\end{ejercicio}
\begin{solucion}
Elegimos conjuntos crecientes $E_n$ de medida finita que cubran $\{f\ne0\}$ y ponemos $f_n=f\ind{E_n\cap\{|f|\le n\}}$. Entonces $f_n\in L^q$, $|f_n|\uparrow|f|$, y los funcionales $T_n(g)=\int f_ng$ tienen norma $\norm{f_n}_q$ por la proposición \ref{v2:dual-isometria}. Para cada $g\in L^p$,
$|T_n(g)|\le\int|fg|<\infty$, uniformemente en $n$. El principio de acotación uniforme, aplicado al Banach $L^p$, da $\sup_n\norm{f_n}_q=M<\infty$. La convergencia monótona implica $\int|f|^q=\lim_n\int|f_n|^q\le M^q$. Así $f\in L^q$.
\end{solucion}
\begin{ejercicio}[Lectura 21, ejercicio 3]\label{ej:v2-l21-3}
Supóngase que $\mu$ es semifinita: cada conjunto de medida infinita contiene uno de medida positiva finita. Si $f$ es medible y $fg\in L^1$ para todo $g\in L^p$, demostrar $f\in L^q$.
\end{ejercicio}
\begin{solucion}
Primero observemos que, en un conjunto $A$ de medida infinita, las medidas de sus subconjuntos de medida finita no pueden estar acotadas. Si su supremo fuera $M<\infty$, una sucesión que se acercara a $M$ tendría unión $B$ de medida $M$, tomando uniones finitas crecientes. El resto $A\setminus B$ tendría medida infinita y, por semifinitud, contendría un $D$ de medida positiva finita. Entonces $B\cup D$ contradice la definición de $M$. Por tanto podemos elegir $B_n\subset A$ de medida finita con $\mu(B_n)>n$; su unión es $\sigma$-finita e infinita.

Para $A_k=\{|f|\ge1/k\}$, si $\mu(A_k)=\infty$, elegimos un $B\subset A_k$ de medida $\sigma$-finita e infinita. La función $f\ind B$ satisface el ejercicio \ref{ej:v2-l21-2}, luego pertenece a $L^q$, pero $\int_B|f|^q\ge k^{-q}\mu(B)=\infty$, contradicción. Así todos los $A_k$ tienen medida finita. Su unión es $\{f\ne0\}$, y el ejercicio anterior se aplica directamente a $f$.
\end{solucion}
\begin{ejercicio}[Lectura 21, ejercicio 4]\label{ej:v2-l21-4}
En $\R$, sea $\Acal$ la $\sigma$-álgebra de conjuntos numerables o con complemento numerable, y sea $\mu(A)=0$ en el primer caso y $\mu(A)=1$ en el segundo. Explicar por qué la petición de alcanzar ambos extremos del ejercicio \ref{ej:v2-l21-1} falla en este espacio para $p\ne2$.
\end{ejercicio}
\begin{solucion}
Toda función medible real es constante fuera de un conjunto numerable. Para verlo, dividimos su recorrido en intervalos de longitud $2^{-n}$. De sus preimágenes, exactamente una puede ser conumerable y al menos una debe serlo, pues una unión numerable de conjuntos numerables no cubre $\R$. Fuera de la unión numerable de las preimágenes descartadas, los valores de la función están a distancia a lo sumo $2^{-n}$ para todo $n$, luego son iguales. Para funciones complejas se aplica el argumento a sus partes real e imaginaria. Por tanto cada $L^p$ es isométrico a $\F$ y el cociente del ejercicio es siempre uno, por el paralelogramo escalar. Las desigualdades siguen siendo verdaderas, pero ninguno de sus extremos estrictos se alcanza si $p\ne2$. Dos conjuntos de medida positiva pueden coincidir o intersectarse en medida uno; su mera existencia no sustituye la disjunción.
\end{solucion}
\begin{ejercicio}[Lectura 21, ejercicio 5]\label{ej:v2-l21-5}
Una sucesión $v_n$ converge débilmente a $v$ si $\varphi(v_n)\to\varphi(v)$ para todo $\varphi\in V^*$. Un operador continuo $T:V\to W$ es completamente continuo si transforma convergencia débil en convergencia en norma. Para espacios de Banach $V,W$, demostrar: (i) todo operador compacto es completamente continuo; (ii) si $V$ es reflexivo, todo operador completamente continuo es compacto.
\end{ejercicio}
\begin{solucion}
(i) Una sucesión débilmente convergente es acotada: se aplica acotación uniforme a los funcionales $J(v_n)$ sobre $V^*$ y se usa $\norm{J(v_n)}=\norm{v_n}$, consecuencia de Hahn--Banach. Si $T$ es compacto, toda subsucesión de $Tv_n$ admite una subsucesión convergente en norma. Su límite debe ser $Tv$, porque para $\psi\in W^*$, $\psi(Tv_n)=(T^*\psi)(v_n)\to\psi(Tv)$, y los funcionales separan puntos por Hahn--Banach. Si $Tv_n$ no convergiera en norma a $Tv$, una subsucesión quedaría a distancia al menos $\eps>0$ y su subsucesión convergente produciría una contradicción.

(ii) Justificamos el hecho de extracción que se necesita. En un Banach reflexivo toda sucesión acotada tiene una subsucesión débilmente convergente. Sea $Y$ el subespacio cerrado generado por los términos de la sucesión; es separable. También es reflexivo: dado $F\in Y^{**}$, la composición de $F$ con la restricción $V^*\to Y^*$ pertenece a $V^{**}$, luego es evaluación en algún $v\in V$. Todo funcional que se anula en $Y$ se anula en $v$; la separación de un punto de un subespacio cerrado mediante Hahn--Banach implica $v\in Y$. Así $F$ es evaluación en ese $v$.

El dual $Y^*$ es separable. En efecto, elijamos $y_j$ densos en la esfera de $Y$ y funcionales $\varphi_j$ de norma uno con $\varphi_j(y_j)=1$. La familia es normante: si $y\ne0$, aproximar $y/\norm y$ por $y_j$ muestra $\sup_j|\varphi_j(y)|=\norm y$. Si el espacio cerrado generado por los $\varphi_j$ no fuera $Y^*$, Hahn--Banach daría un funcional no nulo de $Y^{**}$ que se anula sobre todos ellos; por reflexividad sería evaluación en un $y\ne0$, contradiciendo que la familia sea normante. Las combinaciones racionales de los $\varphi_j$ forman, pues, un conjunto denso numerable de $Y^*$.

Una extracción diagonal hace converger las evaluaciones de una sucesión acotada sobre ese conjunto numerable. La cota uniforme de las normas extiende la convergencia a todo $\varphi\in Y^*$ y define un funcional continuo $F(\varphi)=\lim\varphi(v_{n_j})$. Reflexividad da $F=J(y)$, es decir, convergencia débil a $y$. La restricción de cualquier funcional de $V^*$ muestra que también es débil en $V$.

Aplicamos este hecho a cualquier sucesión de la bola unitaria de $V$. Su subsucesión débilmente convergente tiene, por completa continuidad, imagen convergente en norma. Así la imagen de la bola es relativamente secuencialmente compacta; en un espacio métrico completo eso equivale a compacidad de su cerradura. Por tanto $T$ es compacto.

La lectura menciona además un resultado atribuido a Odell acerca de espacios que no contienen $\ell^1$. Esa afirmación pertenece a la teoría de sucesiones débiles; no se usa como hipótesis ni como paso de estas dos demostraciones.
\end{solucion}
\begin{ejercicio}[Lectura 21, ejercicio 6]\label{ej:v2-l21-6}
Sean $1<p<\infty$, $q=p/(p-1)$ y $\mu$ $\sigma$-finita. Sea $k:X\times X\to\F$ medible para $\Acal\otimes\Acal$, con
\[
 \alpha=\int_X\left(\int_X|k(x,y)|^p\,d\mu(x)\right)^{q/p}d\mu(y)<\infty.
\]
Demostrar que $Kf(x)=\int_Xk(x,y)f(y)\,d\mu(y)$ define, casi en todas partes, un operador compacto en $L^p$, con $\norm K\le\alpha^{1/q}$.
\end{ejercicio}
\begin{solucion}
Escribamos $a(y)=\norm{k(\cdot,y)}_p$, de modo que $\norm a_q=\alpha^{1/q}$. La desigualdad integral de Minkowski da
\[
 \norm{Kf}_p\le\int_X a(y)|f(y)|\,d\mu(y)\le\norm a_q\norm f_p. \tag{*}
\]
Aquí también hay que justificar existencia. Para $H(x)=\int|k(x,y)f(y)|\,d\mu(y)$, Tonelli y Hölder implican $\int Hh\le\norm a_q\norm f_p\norm h_q$ para todo $h\ge0$ de $L^q$. Aplicamos esta desigualdad a las funciones que realizan la norma de $H_m=\min\{H,m\}\ind{E_m}$, donde $E_m\uparrow X$ y $\mu(E_m)<\infty$. Obtenemos $\norm{H_m}_p\le\norm a_q\norm f_p$; por convergencia monótona $H\in L^p$ con la misma cota. Así la integral definitoria es absolutamente convergente casi en todas partes y (*) queda probada. También queda claro que cambiar $f$ en un conjunto nulo no cambia $Kf$ en $L^p$.

Demostraremos compacidad aproximando el núcleo. La norma mixta
\[
 \norm h_{q,p}=\left(\int\norm{h(\cdot,y)}_p^q\,d\mu(y)\right)^{1/q}
\]
controla, por (*), la norma del operador asociado. Los núcleos $k_m=k\ind{E_m\times E_m}\ind{\{|k|\le m\}}$ convergen a $k$ en esta norma: primero, para casi todo $y$, convergencia dominada en $x$ da convergencia de las normas $L^p$ del error; después dominada en $y$, con cota $a(y)^q$, da la afirmación.

En un rectángulo de medida finita, todo núcleo acotado puede aproximarse en $L^r$, $r\ge\max\{p,q\}$, por combinaciones finitas de indicadores de rectángulos medibles. Explicamos el punto sobre conjuntos: la familia de conjuntos aproximables en medida por el álgebra de uniones finitas de rectángulos contiene esa álgebra, es estable por complementos relativos y por uniones numerables. Para esta última propiedad se aproxima primero una unión finita, pues la medida total es finita, y después cada uno de sus miembros. Por tanto contiene la $\sigma$-álgebra producto. Aproximar niveles de una función simple da la densidad en $L^r$. Además, si $h$ está soportado en $E\times F$ de medidas finitas positivas, dos aplicaciones de Hölder dan
\[
 \norm h_{q,p}\le\mu(E)^{1/p-1/r}\mu(F)^{1/q-1/r}\norm h_{L^r(E\times F)}.
\]
Los casos de medida cero son triviales. Así obtenemos núcleos $s_j(x,y)=\sum_{i=1}^{m_j}c_i\ind{A_i}(x)\ind{B_i}(y)$ que convergen a $k$ en norma mixta, con $A_i,B_i$ de medida finita. Su operador es
\[
 S_jf=\sum_i c_i\left(\int_{B_i}f\,d\mu\right)\ind{A_i},
\]
y tiene rango finito. Por (*), $\norm{K-S_j}\le\norm{k-s_j}_{q,p}\to0$. Un límite en norma de operadores compactos es compacto: una red finita para una aproximación, junto con el error uniforme sobre la bola unitaria, da acotación total de la imagen; la completitud del codominio da compacidad de la cerradura. Esto completa tanto la definición como la compacidad de $K$.
\end{solucion}
\section{Geometría y representación}
La convexidad cuantitativa de Clarkson obliga a que una sucesión casi maximizante sea de Cauchy. La diferenciación de la norma identifica después el funcional que apoya a la esfera en su límite. El argumento explica por qué el intervalo $1<p<\infty$ tiene una dualidad especialmente regular y por qué los extremos $p=1,\infty$ requieren estudiar propiedades adicionales de la medida.
